\paragraph{Identity of the section and composition of its observables.}
\label{el-ampl-articulacion}
The structural identity of the electron is introduced through the central
section and its transport operations. Inversion of the APP chart
selects the cell $(5,5)$ by its fixed-point property. The two arithmetic
sheets allow the deformations preserving the sum and the product
residue to be studied. The result distinguishes the two orientations
$(8,2)$ and $(2,8)$ and preserves a difference of one quotient unit.
This is the first level of information: a cell, an involution, a
preserved fibre, and a minimal integer defect. No mass evaluation
is required to state or prove it.

The passage to arithmetic projections introduces a second level. The
sum and product fibres determine conditional expectations on the
same finite space. Their commutator measures dependence on the
order of these two projections. The norm $\sqrt3/4$ follows from the spectrum
of their pairing; the block attaining this norm also preserves
the quadratic relations proved below. Thus, the scalar
prefactor and the local algebra have a common provenance, but
contain different information. The norm retains the magnitude of the
commutator, while the matrices preserve its orientation and its
composition law.

Areal--volumetric transfer incorporates the regional coordinates
at a third level. The mode calendar determines the incidence
matrix and its quadratic action. The closure mark produces the ratio
$\pi/\varphi^2$; interchanging the two arguments provides the
orientation $\varphi/\pi^2$ used in the electronic exponential
character. The two expressions are related by an explicit
involution. This relationship explains why a regional coordinate may
appear in the electronic equation with a power or in a position
different from that in its first reading: the operation applied
is part of the mathematical datum that is preserved.

The cubic alpha term belongs to the regional complement of a marked
return. Counting its positions yields $27-5=22$, while the
three sheets of each of the five regions yield $3\cdot5=15$.
The proofs preserve the domain of these counts and the
orientation determining their signs. The basal formula therefore
composes an incompatibility norm, a transfer character,
a position deficit, and a torsional memory.
Writing it on a single line expresses the composition of these
operations; explaining it requires reconstructing them in that order.

The transition from the basal coordinate to the complete coordinate preserves
the same central section. The return factor uses the history
records, and the two stated readers recover the triple
$(80,54,6)$ along that trajectory. Multiplication by
$R_{\mathrm{act}}^{80-54\alpha+6\Delta_4}$ modifies the evaluation,
but keeps its origin and antecedents identified. The basal
value constitutes an intermediate stage of the composition; the complete
value also includes this return contribution. The physical
comparison must use the stage specified by the formula being evaluated.

The action scale intervenes subsequently with a different function. A ratio of
two action sections is dimensionless and cancels their common scale.
The absolute realization instead uses an action basis, a
temporal reader, and the map to energy or mass. This separation allows
the construction of the scale and the identity of the electronic section
to be preserved simultaneously. The chronological quantum provides a
scale reference; the central section determines the structural factor
expressed relative to it. The dimensional equations of this
chapter keep that composition explicit.

The electron's relevance to the overall argument arises from this
convergence of operations. The same arithmetic that produced the
regions now determines a central cell, its deformations, and its
projections. The fine-structure constant, constructed earlier, acts
in this composition as a closure coordinate. Return information
determines the complete evaluation, and dimensional realization
allows subsequent comparison. Explaining the electron from its structure
consists in preserving this succession of objects, maps, and values at every
step of the proof.

